@@M0@@
and it is an isomorphism on each component, since @@M1@@ was assumed strictly perfect, whence ii).

\sgaStableTarget{sga6.e06.sec.2.7}2.7. On @@M2@@, the functor which associates to a graded @@M3@@-Module its @@M4@@-th component is an exact functor with values in @@M5@@; it therefore defines a functor:
\sgaStableTarget{sga6.e06.eq.2-7-1}@@M6@@
For every @@M7@@, we can consider the functor @@M8@@ defined by tensoring with @@M9@@. Homomorphism (1.2.1) defines a morphism of functors on @@M10@@, whence a morphism for the corresponding functors between derived categories:
\sgaStableTarget{sga6.e06.eq.2-7-2}@@M11@@
Let @@M12@@ be a @@M13@@-perfect complex. We shall show that there exists a sequence of distinguished triangles:
@@M14@@
where @@M15@@ is acyclic, and the @@M16@@ are of the form @@M17@@, the @@M18@@ being perfect complexes in the ordinary sense on @@M19@@.

Since @@M20@@ is @@M21@@-perfect, and @@M22@@ is quasi-compact, there exists a finite covering of @@M23@@, and an integer @@M24@@ such that, on every open of this covering, @@M25@@ is isomorphic to a bounded complex whose terms are locally free of finite type over @@M26@@, say @@M27@@, such that @@M28@@ for @@M29@@.

